% chktex-file 0
\documentclass{ctexart}
\usepackage{graphicx} % Required for inserting images
\usepackage[margin=2.5cm]{geometry}
\usepackage{booktabs}
\usepackage{float}
\usepackage{soulutf8}
\usepackage{enumitem}
\usepackage{pifont}
\usepackage{pgfplots}
\usepackage{longtable}
\pgfplotsset{compat=newest}
\usepackage{tikz}
\usepackage{xcolor} % 用于颜色设置
\usepackage{listings} % 用于代码块
\usepackage{fontspec} % 关键：用于使用系统字体 (XeLaTeX/LuaLaTeX)
% Linux 环境中没有 Windows 宋体时，使用可用的中文衬线字体作为宋体替代
\setCJKmainfont{Noto Serif CJK SC}
\usepackage{authblk}
\usepackage[
    colorlinks=true,      % 使用颜色而非方框
    linkcolor=blue,       % 内部链接颜色
    citecolor=green,      % 引用链接颜色
    urlcolor=red,         % URL 链接颜色
    bookmarks=true,       % 生成 PDF 书签
    bookmarksopen=true,   % 默认展开书签
    pdfstartview=FitH     % 页面宽度适应窗口
]{hyperref}

\sethlcolor{yellow}%设置所需高亮颜色
\usepackage[most]{tcolorbox}
\tcbuselibrary{listings,skins,breakable}
% 页眉页脚
\usepackage{fancyhdr}
\pagestyle{fancy}

% 定义草绿色
\definecolor{grassgreen}{RGB}{124,179,66}
% 补全缺失的浅蓝色定义（防止编译报错）
\definecolor{lightblue}{RGB}{30,144,255} 

% 清空默认
\fancyhf{}

% 页眉
\fancyhead[L]{\textcolor{lightblue}{\kaishu Environmental Ethics}}
\fancyhead[R]{\textcolor{lightblue}{\textit{2026.9-2027.1}}}

% 页脚页码居中
\fancyfoot[C]{\textcolor{lightblue}{\thepage}}

% 页眉、页脚线颜色
\renewcommand{\headrulewidth}{1pt}
\renewcommand{\footrulewidth}{1pt}
\renewcommand{\headrule}{\color{blue!30!white}\hrule height 1pt width\headwidth}
\renewcommand{\footrule}{\color{blue!30!white}\hrule height 1pt width\headwidth}
\newtcolorbox{fancybox}[2][blue]{ % 可选颜色参数
  enhanced,
  breakable,
    colback=#1!3!white,
    colframe=#1!18!white,
    boxrule=0.8pt,
    arc=6pt,
    outer arc=6pt,
  left=10pt,right=10pt,top=12pt,bottom=10pt,
  title={\textbf{#2}},
    fonttitle=\bfseries\color{#1!65!black},
  attach boxed title to top left={yshift=-2mm,xshift=5mm},
  boxed title style={
        colback=#1!8!white,
        colframe=#1!18!white,
        boxrule=0.6pt,
        arc=4pt,
  },
    shadow={1.5pt}{-1.5pt}{0pt}{black!10},
}

% 定义一套适合C++的配色方案
\usepackage{tabularx}
\newfontfamily{\consolasfont}{Ubuntu}

\newcommand{\bluecode}[1]{{\color{blue}\consolasfont #1}}
\lstdefinestyle{MyCppStyle}{
    language = C++,
    % 确保这里使用 \consolasfont
    basicstyle = \consolasfont\small, 
    keywordstyle = \color{blue!70!black}, 
    commentstyle = \color{green!50!black}, 
    stringstyle = \color{red!60!black}, 
    numberstyle = \tiny\color{gray}, 
    numbers = left, 
    frame = single, 
    rulecolor = \color{lightgray}, 
    breaklines = true, 
    backgroundcolor = \color{white}, 
    tabsize = 4, 
    showstringspaces = false, 
    xleftmargin = 2em, 
    escapeinside=``,
}
% --- lstdefinestyle HTML 风格定义 ---
\lstdefinestyle{mhtml}{
    language = HTML,
    basicstyle = \consolasfont\small, 
    tagstyle = \color{blue!70!black}\bfseries,
    keywordstyle = \color{purple!80!black},
    commentstyle = \color{green!50!black}\itshape, 
    stringstyle = \color{red!60!black}, 
    numberstyle = \tiny\color{gray}, 
    numbers = left, 
    frame = single, 
    rulecolor = \color{lightgray}, 
    breaklines = true, 
    backgroundcolor = \color{white}, 
    tabsize = 4, 
    showstringspaces = false, 
    xleftmargin = 2em, 
    escapeinside=``,
    morekeywords={als, doctype, html, head, title, meta, link, body, div, span, h1, h2, h3, h4, h5, h6, p, a, img, ul, ol, li, table, tr, th, td, pre, code, script, style},
    morekeywords=[2]{class, id, href, src, rel, type, charset, lang, name, content, style},
    keywordstyle=[2]\color{purple!80!black},
}
\usepackage{hyperref}
\hypersetup{
    colorlinks=true,      
    linkcolor=blue,       
    filecolor=magenta,    
    urlcolor=cyan,        
    pdftitle={我的文档标题}, 
}

% ==================== 正文部分 ====================
\begin{document}

% 开启无页眉页脚模式（针对封面和目录）
\pagestyle{empty}

% ------------------ 美化封面页 ------------------
\begin{titlepage}
    \centering
    % 顶部 TikZ 几何几何装饰条（全宽 bleed 效果）
    \begin{tikzpicture}[remember picture, overlay]
        \fill[blue!50!white] (current page.north west) rectangle ([yshift=-2cm]current page.north east);
        \fill[grassgreen] (current page.north west) ++(0,-2cm) rectangle ([yshift=-2.2cm]current page.north east);
    \end{tikzpicture}
    
    \vspace*{3.5cm}
    
    % 大标题
    {\fontsize{30pt}{36pt}\selectfont\bfseries\color{blue!80!black} 环境伦理学}\\[0.6em]
    % 副标题或英文标识
    {\fontsize{14pt}{18pt}\selectfont\color{gray}\textsf{Environmental Ethics}}\\[1cm]
    
    % 居中装饰线
    \centering\makebox[4cm]{\color{grassgreen}\rule{5cm}{1.5pt}}
    
    \vfill
    
    % 使用 tcolorbox 包装的精致作者信息框
    \begin{tcolorbox}[
        enhanced,
        width=0.55\textwidth,
        colback=blue!5!white,
        colframe=blue!5!white,
        boxrule=1.2pt,
        arc=6pt,
        sharp corners=downhill, % 别致的对角切角效果
        shadow={2pt}{-2pt}{0pt}{black!20},
        halign=center,
        valign=center,
        top=15pt, bottom=15pt
    ]
        {\large\bfseries\color{black!80} Author：} {\large\kaishu Simson} \\[0.6em]
        {\large\bfseries\color{black!80} Date：} {\large\kaishu September 2026}
    \end{tcolorbox}
    
    \vfill
    \vspace*{2cm}
    
    % 底部标识
    {\color{gray!80}\small\the\year\ \copyright\ All Rights Reserved.}
\end{titlepage}

\pagestyle{fancy}

\newpage
\tableofcontents
\newpage

% 设置从当前页（正文第一页）开始采用阿拉伯数字编码，并自动重置为第 1 页
\pagenumbering{arabic}

\section{概论}

环境伦理学研究环境相关的价值判断和伦理判断。相比于一般应用伦理学，要努力突破人类中心主义视角，对环境相关的伦理问题进行反思。环境伦理学的讨论分以下几个部分：

\begin{enumerate}
    \item 环境问题的历史和现状
    \begin{itemize}
        \item 环境问题的历史
        \item 20世纪中期依赖的环境危机
        \item 中国环境现状
    \end{itemize}
    \item 环境危机的根源
    \begin{itemize}
        \item 从经济，技术层面的反思
        \item 从文化价值层面的反思————对环境危机根源认识的深化
    \end{itemize}
    \item 环境伦理学的性质和起源
    \begin{itemize}
        \item 环境伦理学的学科性质
        \item 几种主要理论和观点
        \item 研究方法和研究思路
    \end{itemize}
    \item 伦理边界的扩展
    \begin{itemize}
        \item 伦理的进化
        \item 人类中心主义的主要形态和根据
        \item 非人类中心主义
        \item 两种立场
    \end{itemize}
    \item 自然的道德地位
    \begin{itemize}
        \item 利益原则与道德可考虑性（费因伯格）
        \item 树木具有法律地位吗（斯通）
    \end{itemize}
    \item 动物解放与动物权力论
    \begin{itemize}
        \item 动物的出境和人类道德
        \item Singer和动物解放：功利主义和道德眼神主义
        \item 动物权利论：义务论，天赋价值和生命主体
        \item 动物解放与素食主义
        \item 动物解放与环境伦理学的矛盾
    \end{itemize}
    \item 生物中心环境伦理学
    \begin{itemize}
        \item 生物中心论
        \item 罗尔斯顿
        \item 敬畏生命的伦理学（史怀哲）
        \item 尊重自然的环境伦理学（Paul Taylor）
    \end{itemize}
    \item 生态中心主义环境伦理学
    \begin{itemize}
        \item 生物中心主义环境伦理学的缺陷
        \item 利奥波德与大地伦理————整体主义伦理学
        \item 对整体主义环境伦理的批评
    \end{itemize}
    \item 深层生态学
    \begin{itemize}
        \item 奈斯：浅层的和深层的生态学
        \item 生态智慧与自我实现
        \item 生态哲学与宗教————精神生态学
    \end{itemize}
    \item 环境正义
    \begin{itemize}
        \item 起源：对非人类中心主义misanthropic倾向的批评
        \item 环境社会学的重要研究：环境污染与社会不公的关联
        \item 消除环境不公的根源：国内的和国际的
    \end{itemize}
    \item 社会生态学
    \begin{itemize}
        \item 马克思主义和生态社会主义
        \item 生态无政府主义————Murray Bucchin
        \item 生态乌托邦
    \end{itemize}
    \item 环境伦理的精神资源
    \begin{itemize}
        \item 中西方传统文化中支持环境伦理的人文，哲学和宗教思想
    \end{itemize}
    \item 走向实践的环境伦理
\end{enumerate}

\section{环境伦理学的发展与环境现状}

\subsection{环境问题}

环境问题的本质是人类与自然在物质和能量的交换上产生的问题。

\subsubsection{环境问题的历史}

\paragraph{环境史}
\begin{enumerate}
    \item 农耕时代：人的文化和社会是自然的循环的一部分。也有环境问题，但是和今天有本质的区别。
    \item 工业时代：利用化石燃料。城市的兴起令人类脱离自然
    \item 二战之后：化学合成物增多，人类或冻对地球的影响越来越大。
\end{enumerate}

在战后，用于武器的技术变化为核电站，DDT变化为农用杀虫剂。20世纪的环境危机体现在物流师年代开始各种公害事件频繁出现，比如：

\begin{itemize}
    \item LA光化学烟雾事件（50年代初）
    \item 日本水俣病（1956）
    \item 伦敦烟雾事件（1952）
\end{itemize}

七十年代以后，环境问题在范围上变得更广，危害人整体的生存和发展。水污染，大气污染，沙漠化和水土流失。关于环境治理的国际责任引发了复杂的环境正义问题的争端。

以上的历史现实催生了现代环境运动，和民权运动等相互结合，有以下几个重要的人物：

\paragraph{现当代环境运动}

\begin{itemize}
    \item 1962年雷切尔卡森《寂静的春天》，表达了对人通过技术来控制自然的批评。这是现代环境运动的先驱。
    
    思想谱系：继承史怀哲（序言）“敬畏生命”的观点，表达了对于人类中心主义的激烈批判。
    \item 1967历史学家霖林恩·怀特在《科学》发文“我们生态危机的历史根源”，指责基督教为生态危机负责。他揭示了基督教发展的历史和西方基督教文明对环境的破坏。
    \item 1968年经济学家提出“公地悲剧”，将环境问题拔高到经济学和社会学高度。（观点：需要靠胁迫性的共识，如法律、义务，形成纳税一类义务）
    
    \textbf{著名口号}：生育权不是一种人权，号召采取全球性政策阻止贫困国家大量生育。
    \item 1971年巴里·康芒纳《封闭的循环————自然，人和技术》，技术的生态缺陷来源于背后的科学方法。
    \item 1972年罗马俱乐部发表第一个研究报告《增长的极限》，语言经济增长不可能无限持续下去，因为石油等自然资源的供给有限，预测世界性灾难即将来临
    
    对策：提出“零增长”对策。（The limit to growth）
    \item 1973英国经济学家舒马赫在《小的是美好的》一书中对于西方经济学的环境后果做了分析和批评。
\end{itemize}


\end{document}